\section{Output Analysis}

\subsection{Design}

Along with the AS-IS model, three TO-BE scenarios are evaluated as described in section 1.2.
As each session lasts approximately 4 hours, we set the simulation time to 6 hours to ensure that every patient could finish their visit. All models were run with 30 replications, which gives a 95\% confidence interval of all replications below 8 minutes. The random seeds used in the probability distributions were consistent for all models, enabling comparison and reducing variance across simulations.

\subsection{Performance measures}

The primary output measure is mean cycle time per patient \parencite{MillerPainTreatmentCenter}, defined as time in minutes from arrival to checkout. This directly affects patient care, and is Dr. Weems' primary focus in optimizing the system. Total patients processed is a secondary measure, to ensure that enhancements made do not affect overall system performance.

\subsection{Confidence intervals}

\autoref{tab:ci} presents mean cycle time and 95\% confidence intervals for all models across 30 replications. The results show a clear improvement of the system for all modifications, with the combined solution giving the best performance. The reduction in confidence interval show not only an improvement in performance, but also a reduction in variability, suggesting a more consistent and reliable process flow.

\begin{table}[h]
\centering
\begin{threeparttable}
\begin{tabular}{lrrrr}
\toprule
\textbf{Model} & \textbf{Mean} & \textbf{Std} & \textbf{CI lower} & \textbf{CI upper} \\
\midrule
AS-IS                   & $73.95$ & $20.81$ & $66.18$ & $81.72$ \\
TO-BE 1: Scheduling     & $65.60$ & $15.34$ & $59.87$ & $71.33$ \\
TO-BE 2: Pre-processing & $60.99$ & $14.80$ & $55.46$ & $66.51$ \\
TO-BE 3: Combined       & $53.95$ & $10.31$ & $50.10$ & $57.80$ \\
\bottomrule
\end{tabular}
\begin{tablenotes}
\small
\item 
\end{tablenotes}
\end{threeparttable}
\caption{Mean cycle time and 95\% confidence intervals across models}
\label{tab:ci}
\end{table}

\subsection{Statistical comparison}

\autoref{tab:ci} shows that none of the TO-BE scenario has overlapping confidence intervals with the AS-IS simulation, suggesting statistically significant differences \parencite[p. 390]{pensumbok}.

This suggestion is confirmed by the paired t-test applied to compare the outputs from each simulation, measuring the average deviation between models \parencite[p. 426-429]{pensumbok}. Since random seeds are shared for all models, performance difference $D_i = \bar{x}_i^{TO-BE} - \bar{x}_i^{AS-IS}$ is a result of enhancements made rather than variance. The null hypothesis $H_0: \mu_D = 0$ is tested at the 5\% significance level with $n - 1 = 29$ degrees of freedom.

The null-hypothesis is rejected for all scenarios, as the t-statistic is well below the critical value of $t_{0.025,29} = -2.045$. 

\begin{table}[H]
\centering
\begin{threeparttable}
\begin{tabular}{lrrrrr}
\toprule
\textbf{Comparison} & \textbf{$\bar{D}$} & \textbf{CI lower} & \textbf{CI upper} & \textbf{t-statistic} & \textbf{p-value} \\
\midrule
Scheduling vs.\ AS-IS      & $-6.71$  & $-8.12$  & $-5.31$ & $-9.80$  & $<0.001$ \\
Pre-processing vs.\ AS-IS  & $-5.47$  & $-6.76$  & $-4.19$ & $-8.71$  & $<0.001$ \\
Combined vs.\ AS-IS        & $-12.11$ & $-14.41$ & $-9.81$ & $-10.76$ & $<0.001$ \\
\bottomrule
\end{tabular}
\end{threeparttable}
\caption{Paired t-test results}
\label{tab:ttest}
\end{table}